% ============================================================
% Section 02 — Identity and Persistence
% Book I: Ontology of Viability
% Chapter 05: Identity and Continuity
%
% Purpose:
%   To distinguish the persistence of an existent from
%   the persistence of its identity and from the continuity
%   of identity across temporal succession.
%
% Conceptual Policy:
%   1. Identity is not reducible to persistence.
%   2. Persistence of an existent does not automatically
%      establish persistence of identity.
%   3. Persistence, continuity, and temporal succession
%      must remain analytically distinct.
%   4. Material, structural, functional, relational, and
%      manifestational preservation are not independently
%      sufficient criteria of identity persistence.
%   5. Neither the termination nor the preservation of
%      identity may be inferred solely from the interruption
%      or continuation of a particular manifestation.
% ============================================================

\section{Identity and Persistence}
\label{sec:identity-and-persistence}

\subsection{The Ontological Distinction Between Identity and Persistence}

Identity and persistence are closely related ontological
notions, but they do not designate the same condition.

Identity concerns the numerical and ontological identity
of an existent: the question of what makes an existent
the existent that it is, without reducing this question
to any particular determination, configuration, or
temporal condition.

Persistence concerns the continuation of an existent
through a specified interval or under specified conditions.
It addresses whether an existent remains in existence
across a domain of temporal or other ordered variation.

The distinction is necessary because the continuation
of an existent and the persistence of its identity are
not automatically equivalent.

An existent may persist through a temporal interval while
undergoing substantial alterations in its determinations,
structure, relations, material constitution, or
manifestations. Conversely, the cessation of a particular
configuration or manifestation does not, by itself,
establish that the identity associated with it has ceased
to exist.

Persistence therefore cannot serve as an immediate
substitute for an ontological account of identity.

The present section examines the relation between these
notions without assuming that persistence necessarily
preserves identity or that the interruption of persistence
necessarily establishes the termination of identity.

\subsection{Persistence of an Existent and Persistence of Identity}

A distinction must be established between the persistence
of an existent and the persistence of its identity.

The persistence of an existent concerns the continuation
of that existent through a specified interval or domain.
The persistence of identity concerns whether the numerical
identity of that existent remains the same across the
relevant transformations or conditions.

These questions are related, but the answer to one does
not automatically determine the answer to the other.

The persistence of an existent may be established in
a particular respect without resolving whether its
identity has remained numerically the same. For example,
the continued presence of a material configuration,
the uninterrupted operation of a process, or the
maintenance of a recognizable structural organization
may establish certain forms of persistence. None of
these facts alone establishes the persistence of identity.

Likewise, a failure to establish the uninterrupted
persistence of a particular manifestation does not
necessarily establish the non-persistence of identity.

The distinction prevents the ontological question from
being settled through the mere observation of temporal
continuation or interruption.

Accordingly, every claim concerning persistence must
specify what is said to persist: the existent, its identity,
a determination, a structure, a relation, a process,
or a manifestation.

\subsection{Persistence and Temporal Continuity}

Persistence and temporal continuity must also be
distinguished.

Persistence refers to the continuation of an existent
through a specified interval. Temporal continuity concerns
the relation between successive temporal conditions
and whether the relevant succession is uninterrupted
or otherwise continuous according to a defined criterion.

An existent may be described as persistent even when
a particular process or manifestation associated with
it is interrupted. Similarly, temporal continuity in
a process does not necessarily establish the persistence
of the identity of the existent involved in that process.

The continuity of a process and the persistence of
an existent therefore belong to different analytical
levels.

A process may continue while its underlying conditions
or participating entities change. An existent may also
remain in existence despite interruptions in particular
activities or manifestations.

It follows that uninterrupted temporal succession
cannot be treated as a universal criterion of identity
persistence.

Nor can the presence of an interruption be treated
as sufficient evidence that identity has ceased.

The ontological significance of an interruption depends
on what has been interrupted, the domain in which the
interruption occurs, and the respect under which
persistence is being examined.

\subsection{Persistence and the Continuity of Manifestation}

The persistence of identity must not be conflated with
the continuity of manifestation.

Manifestations may continue, change, diminish, become
interrupted, or cease under different conditions.
Their continuity concerns the persistence or succession
of particular forms through which an existent is manifested.

Identity, however, cannot be identified with the
continuous presence of any particular manifestation.

The continued existence of a manifestation does not
independently establish that the identity associated
with it has persisted. A manifestation may remain
continuous while the ontological identity of its source
or bearer is transformed or otherwise brought into question.

Conversely, the interruption or absence of a manifestation
does not independently establish that the corresponding
identity has ceased to persist.

This distinction also applies to the capacity to manifest.
The continuity of the capacity to manifest is analytically
different from the continuity of actual manifestations
and from the persistence of identity.

An existent may retain a capacity without continuously
actualizing it. A manifestation may be interrupted
without thereby establishing the disappearance of
the capacity associated with it. Yet the persistence
of such a capacity, considered alone, does not establish
the persistence of identity.

Three questions must therefore remain distinct:

\begin{enumerate}
    \item Does the identity of the existent persist?
    \item Do its manifestations continue?
    \item Does its capacity to manifest continue?
\end{enumerate}

No one of these questions can be answered solely by
answering either of the other two.

\subsection{Persistence Through Material and Structural Change}

The persistence of an existent may be examined in relation
to changes in its material constitution and structural
organization.

Material continuity is one possible respect under which
persistence may be investigated. An existent may retain
some of its material constituents across a specified
interval, undergo partial material replacement, or
experience extensive material transformation.

However, the persistence of material constituents does
not by itself establish the persistence of identity.
Nor does the replacement or alteration of material
constituents automatically establish the emergence
of a numerically different identity.

The same caution applies to structural persistence.

A structure may remain relatively stable while its
components, relations, or material realization change.
Alternatively, a structure may undergo substantial
reorganization while the question of the identity of
the existent remains unresolved.

Structural preservation and structural transformation
are therefore relevant to the investigation of persistence,
but neither can be treated as a universal and sufficient
criterion of identity.

The ontological significance of material and structural
change must be determined in relation to the specific
existent and the relevant conditions of identity.

\subsection{Persistence Through Functional and Relational Change}

Persistence may also be investigated through functional
and relational conditions.

An existent may retain a function while undergoing
changes in its internal organization or external relations.
It may also acquire new functions, lose previous functions,
or enter into different relational configurations.

Functional continuity does not independently establish
identity persistence. The same function may be realized
by numerically different existents, and a single existent
may undergo functional transformation without this fact
alone establishing a change of identity.

Likewise, relational continuity cannot be treated as
a sufficient criterion of identity persistence.

An existent may remain related to the same entities
while undergoing transformations in other respects.
It may also enter into new relations or lose existing
relations without the ontological question of identity
being thereby settled.

Functions and relations may contribute to the conditions
under which persistence is investigated. Their significance,
however, must be established rather than assumed.

The persistence of a role, function, or relational position
must therefore be distinguished from the persistence
of the identity of the existent occupying that role
or position.

\subsection{Interruption, Absence, and the Problem of Persistence}

The problem of persistence becomes particularly difficult
when an existent or one of its manifestations undergoes
interruption.

An interruption may concern a process, an activity,
a manifestation, a relation, or a particular condition
of existence. These forms of interruption must not
be treated as ontologically equivalent.

The cessation of a process does not necessarily establish
the cessation of the existent associated with it.
The absence of a manifestation does not necessarily
establish the absence of its source or bearer.
Likewise, the discontinuation of a particular relation
does not necessarily establish the termination of
the identity of the entities involved.

The possibility of interruption therefore requires
a distinction between the interruption of a particular
mode of persistence and the termination of identity.

A further question concerns whether an existent can
persist through an interval during which some of its
manifestations or processes are absent.

This possibility cannot be rejected merely because
the relevant manifestation is not continuously present.
Nor can it be affirmed without examining the ontological
conditions of the existent concerned.

The relation between interruption and persistence
must therefore remain subject to the domain and respect
under examination.

\subsection{Temporary Cessation and the Possibility of Resumption}

The possibility of resumption introduces a further
distinction between uninterrupted persistence and
persistence across an interruption.

A process or manifestation may cease and subsequently
resume. The resumption may involve the continuation
of a previously existing process, the emergence of
a new process, or a relation between distinct processes.
The mere similarity between the conditions before
and after the interruption does not settle which of
these possibilities obtains.

Likewise, the resumption of a manifestation does not
by itself establish that the identity associated with
the earlier manifestation has persisted throughout
the interval.

Nor does the interruption necessarily establish that
the identity has ceased and subsequently been recreated.

The ontological relation between the conditions before
and after an interruption requires independent examination.
It cannot be resolved solely through the temporal
distance between them, the similarity of their
manifestations, or the restoration of a previous
configuration.

The possibility of persistence across interruption
must therefore be distinguished from the uninterrupted
continuity of a process or manifestation.

\subsection{Persistence and Causal Continuity}

Persistence must also be distinguished from causal
continuity.

A causal sequence may connect successive conditions,
events, or configurations. Such a sequence may explain
how one condition produces or contributes to another.
However, causal connection alone does not establish
that the numerically same identity persists across
the sequence.

Causal continuity may connect distinct existents,
and the persistence of an existent cannot be inferred
solely from the existence of a causal relation between
successive conditions.

Conversely, the absence of an established causal
sequence does not, without further argument, settle
the question of identity persistence.

Causal relations may provide relevant conditions
for investigating persistence, but they cannot replace
the ontological analysis of identity.

It is therefore necessary to distinguish the continuity
of causal production from the persistence of the identity
of the entities or processes involved.

\subsection{Persistence and the Conditions of Identity}

The persistence of identity cannot be established
through the isolated preservation of a single feature
without an adequate account of its ontological relevance.

Material constitution, structural organization,
functional capacity, relational configuration,
manifestational continuity, and causal connection
may all be relevant to the investigation. Nevertheless,
none of them can be assumed in advance to constitute
a universally sufficient condition of identity persistence.

The same feature may have different significance
in different domains or for different kinds of existents.
A condition that is essential to the persistence of
one existent may be non-essential to another.

Furthermore, the preservation of a feature in one
respect does not establish the preservation of identity
in every respect.

The investigation must therefore distinguish between:

\begin{itemize}
    \item conditions that accompany persistence;
    \item conditions that contribute to persistence;
    \item conditions that may be necessary for persistence
    in a specified domain;
    \item conditions that may be sufficient for persistence
    under explicitly defined circumstances;
    \item conditions that are merely correlated with
    the persistence of identity.
\end{itemize}

These categories cannot be treated as interchangeable.

The ontological status of each proposed condition must
be established through argument. No criterion of identity
persistence should be adopted merely because it describes
a frequently observed form of continuity.

\subsection{The Non-Equivalence of Persistence and Identity Preservation}

The preceding distinctions establish that persistence
and identity preservation are related but non-equivalent
notions.

Persistence may describe the continuation of an existent,
a process, a manifestation, a structure, or a relation.
Identity preservation concerns the continued numerical
identity of the existent under examination.

The persistence of one of these conditions does not
automatically entail the persistence of the others.

This non-equivalence has two important consequences.

First, the continued existence of something through
a specified interval cannot, without further argument,
be treated as proof that its identity has remained
numerically the same.

Second, the interruption or cessation of a particular
form of persistence cannot, without further argument,
be treated as proof that identity has ceased.

The relation between persistence and identity preservation
must therefore be established through an account of
the relevant ontological conditions, rather than through
the unqualified use of the term persistence.

\subsection{The Subject and Respect of Persistence}

Every ontological claim concerning persistence must
identify its subject and its respect.

The subject may be an existent, an identity, a process,
a structure, a material configuration, a relation,
a capacity, or a manifestation.

The respect may concern temporal duration, material
constitution, structural organization, functional
continuity, relational configuration, causal succession,
or another explicitly defined domain.

Without these distinctions, claims of persistence
remain ambiguous.

For example, the statement that an existent persists
may refer to its continued existence, the preservation
of its numerical identity, the continuity of its
manifestations, or the uninterrupted operation of
a process. These claims are not equivalent and require
different forms of ontological justification.

The subject and respect of persistence must therefore
be specified before any conclusion concerning identity
preservation can be drawn.

\subsection{Methodological Requirements}

The investigation of identity and persistence must
observe several methodological requirements.

First, persistence must not be used as a substitute
definition of identity.

Second, the persistence of an existent must be distinguished
from the persistence of its identity.

Third, temporal continuity, causal continuity, and
continuity of manifestation must be examined separately.

Fourth, material, structural, functional, and relational
preservation must not be treated as universally sufficient
criteria of identity persistence.

Fifth, interruption, absence, and resumption must be
analyzed according to the particular domain and respect
involved.

Sixth, neither the continuation nor the termination
of identity may be inferred solely from the continuation
or cessation of a particular process or manifestation.

Finally, the ontological conditions of identity persistence
must remain open to investigation rather than being
determined in advance by a preferred model of continuity.

\subsection{Scope of the Present Section}

This section has distinguished identity from persistence
and examined the different senses in which an existent,
its identity, its manifestations, and its associated
processes may continue or cease.

The central conclusion is that persistence does not
provide an independent or sufficient explanation of
identity preservation.

The persistence of an existent, the persistence of
its identity, the continuity of its manifestations,
and the continuity of its capacity to manifest are
distinct ontological questions. Their relations must
be established through further analysis.

The next section examines the relation between invariance
and continuity. It investigates whether the preservation
of invariant features can contribute to an account
of identity continuity, and under what conditions
such preservation may or may not have ontological
significance.

% ============================================================
% End of Section 02
% ============================================================